\section{The three matter descriptions and their symmetry limits}
\subsection{Model 1: spinless perfect fluid}
Set the spin current and torsion to zero. A conserved-current action of the standard Brown type \cite{Brown} suffices to make the constitutive content explicit:
\begin{equation}\label{eq:fluidaction}
 I_{\rm pf}=\int\left[-\sqrt{-g}\,\eps(n,\chi)+J^\mu(\partial_\mu\varphi+\beta\partial_\mu\chi+\beta_A\partial_\mu\alpha^A)\right]\dd^4x,
 \quad n=\frac{\sqrt{-g_{\mu\nu}J^\mu J^\nu}}{\sqrt{-g}}.
\end{equation}
$J$ is a vector density; its potential enforces particle conservation, and $\beta$ enforces $u\cdot\nabla\chi=0$. Material labels allow general fluid variations; they introduce no additional stress when eliminated. At fixed current density $J$, $\delta n=\tfrac n2(g_{\mu\nu}+u_\mu u_\nu)\delta g^{\mu\nu}$ and $\delta\sqrt{-g}=-\tfrac12\sqrt{-g}g_{\mu\nu}\delta g^{\mu\nu}$. Substituting these variations into the material energy term gives the perfect-fluid metric source and $P=n\eps_{,n}-\eps$. Varying $J$ and the potentials gives
\begin{equation}
 \nabla_\mu(nu^\mu)=0,\quad \Dt\chi=0,\quad
 T_{\mu\nu}=(\eps+P)u_\mu u_\nu+Pg_{\mu\nu},\quad
 P=n\eps_{,n}-\eps.
\end{equation}
For $\eps=\rho=C(\chi)n^{1+w}$, $P=p=w\rho$ and $\tau=0$. Equations \eqref{eq:energy}--\eqref{eq:momentum} reduce to
\begin{equation}\label{eq:PFconservation}
 \Dt\rho=-(1+w)\rho\Theta,\qquad p'=-(\rho+p)\Phi'.
\end{equation}
For $\rho>0$, $0<w\leq1$ the fluid acoustic speed is $c_s^2=w$, real and causal; the full gravitational gauge and boundary problem must still be specified. Dust $w=0$ has coincident pressureless characteristics and requires a pre-caustic Lagrangian/dust formulation, not a strictly hyperbolic acoustic solver. For $w<0$ or $w>1$ the standard causal Euler interpretation fails; no such branch is authorized by this audit.

\subsection{Model 2: unpolarized, averaged Weyssenhoff fluid}
The microscopic source obeys a convective Frenkel condition, while
\begin{equation}
 \langle S^H_{\mu\nu}\rangle=0,\qquad
 S_H^2=\tfrac12\langle S^H_{\mu\nu}S_H^{\mu\nu}\rangle\ne0.
\end{equation}
Under local rotationally isotropic statistics in the rest frame,
\begin{align}
 \langle S^H_{ij}S^H_{kl}\rangle&=\frac{S_H^2}{3}(\delta_{ik}\delta_{jl}-\delta_{il}\delta_{jk}),\\
 \langle S^H_{i}{}^k S^H_{jk}\rangle&=\frac23 S_H^2\delta_{ij}.
\end{align}
We must first eliminate microscopic torsion and then average: $\langle K^2\rangle$ is not $\langle K\rangle^2$. Setting the mean torsion to zero at the start would erase the effect being modelled. The source used by the existing laboratory is
\begin{equation}\label{eq:avgsource}
 \eps=\rho-q,\quad P=w\rho-q,\quad q=\frac\kap4 S_H^2,\quad Q=0,\quad P_r=P_\perp=P.
\end{equation}
This follows the selected normalization in \cite{Hashemi}; the following \emph{additional} hypotheses are required for an inhomogeneous local use:
\begin{enumerate}
 \item Averaging is over a length small compared with resolved metric/matter variation, large compared with microscopic spin scales.
 \item Derivatives and averaging commute at retained order; correlations of fluctuating velocity, connection, acceleration and shear with spin are discarded or shown negligible. Vanishing mean spin alone is insufficient.
 \item The second moment is locally isotropic, and its statistics and spin per particle are transported without torque, relaxation, spin exchange or particle production.
 \item Metric backreaction/averaging commutators and unresolved anisotropic moments are absent at the approximation order.
\end{enumerate}
These hypotheses are constitutive and coarse-graining assumptions, not established by inserting a radial dependence in a homogeneous formula. Under them, a precise local effective action is Eq.~\eqref{eq:fluidaction} with
\begin{equation}\label{eq:effectiveEOS}
 \eps(n,\chi)=C(\chi)n^{1+w}-D(\chi)n^2,
 \qquad \rho=C(\chi)n^{1+w},\quad q=D(\chi)n^2.
\end{equation}
$D\geq0$ encodes conserved spin variance, not an arbitrary time-dependent repulsion. Variation yields $P=w\rho-q$ and
\begin{align}
 \Dt\rho&=-(1+w)\rho\Theta,&\Dt q&=-2q\Theta,&
 \Phi'&=-\frac{w\rho'-q'}{(1+w)\rho-2q}.\label{eq:avgtransport}
\end{align}
Equivalently $q/\rho^{2/(1+w)}$ is advected. Either $\rho,q$ or $n,\chi$ can be used, not both pairs as independent fluid degrees of freedom. Separate ordinary/spin conservation is part of the model; effective conservation alone would not determine two matter functions. The current action justifies the \emph{effective metric theory} under these hypotheses; it is not a substitute derivation of every microscopic Weyssenhoff action.

\subsubsection{Closure is not hyperbolicity}
Let $h=\eps+P=(1+w)\rho-2q$. At fixed advected composition, the acoustic principal symbol has
\begin{equation}\label{eq:sound}
 c_{\rm eff}^2=\left.\frac{\dd P}{\dd\eps}\right|_\chi
 =\frac{w(1+w)\rho-2q}{(1+w)\rho-2q}.
\end{equation}
Derive this by differentiating Eq.~\eqref{eq:effectiveEOS} with respect to $\ln n$. In a local inertial frame, the frozen-composition principal equations are
\begin{equation}
 \partial_t\delta\eps+h\partial_x\delta v=0,
 \qquad h\partial_t\delta v+\partial_x\delta P=0,
 \qquad \partial_t^2\delta\eps-c_{\rm eff}^2\partial_x^2\delta\eps=0.
\end{equation}
For $0\leq w<1$, the branch $h>0$ has $c_{\rm eff}^2<0$ whenever
\begin{equation}
 \frac{w(1+w)}2<\frac q\rho<\frac{1+w}2.
\end{equation}
In particular, ordinary dust with any $0<q<\rho/2$ has negative acoustic speed squared. Fourier modes grow as $\exp(|k|\sqrt{-c_{\rm eff}^2}t)$, without a wavelength-independent estimate: a generic smooth radial Euler initial-value problem is not well posed on this branch. At $h=0$ the usual Euler inversion degenerates. On the $h<0$ branch $c_{\rm eff}^2>1$ for $w<1$; this is a failure of a causal effective-fluid interpretation, not a theorem that the underlying EC action has propagating ghosts.

A flat homogeneous bounce has $q=\rho$, hence $c_{\rm eff}^2=w+2$. Its trajectory crosses the $h=0$ ratio $(1+w)/2$ before bouncing when these ratios are attained continuously. The spatially homogeneous ODE has no acoustic gradients and remains a consistent symmetry-restricted calculation; this observation does \emph{not} invalidate its existing numerical benchmark. It does prevent calling its unrestricted inhomogeneous continuation a verified causal PDE. For a positive-pressure low-spin regime, only states with $h>0$ and $0<c_{\rm eff}^2\leq1$ are candidates for a causal Euler implementation; approaching $c_{\rm eff}^2=0$ also requires a characteristic-degeneracy analysis. Artificial viscosity, a guessed pressure or changing spin transport is not an authorized cure.

\subsection{Model 3: selected full Weyssenhoff field formulation}
The phrase ``full Weyssenhoff'' does not select a unique averaging scheme or microscopic material action. We use the \emph{published field equations} \cite{BHL}, Eqs. (19)--(27), based on \cite{OK}. This is an allowed identification of an established formulation, not a newly manufactured action. Its full constrained matter-action derivation is an explicitly open reproducibility task.

Keep the BHL spin variable $b_{\mu\nu}$ for linear terms and use their Appendix B transformation to our metric signature: $g_B=-g$, $u_B^\mu=u^\mu$, $(u_B)_\mu=-u_\mu$, $(b_B)_{\mu\nu}=b_{\mu\nu}$. The resulting symmetric effective metric source and spin equation are
\begin{align}
 T_{\mu\nu}^{B}&=(\rho+p-2\kap b^2)u_\mu u_\nu+(p-\kap b^2)g_{\mu\nu}+\mathcal D_{\mu\nu},\label{eq:fullstress}\\
 \mathcal D_{\mu\nu}&=-2(g^{\alpha\beta}-u^\alpha u^\beta)\nabla_\alpha[u_{(\mu}b_{\nu)\beta}],\label{eq:Dspin}\\
 h_\mu{}^\alpha h_\nu{}^\beta\nabla_\lambda(u^\lambda b_{\alpha\beta})&=0,
 \qquad b_{\mu\nu}u^\nu=0,\quad h_\mu{}^\nu=\delta_\mu^\nu+u_\mu u^\nu.\label{eq:fulltransport}
\end{align}
The tensor in parentheses in Eq.~\eqref{eq:Dspin} is \emph{not} the spatial projector; replacing its minus sign by plus without transforming the source equation changes the model. We use $b^2=\tfrac12 b_{\mu\nu}b^{\mu\nu}$. The spin current in the BHL source convention is $u^\lambda b_{\mu\nu}$, their torsion is $Q_B=\kap u b$, and the operational source for Eq.~\eqref{eq:cartanform} is $-2u b$. This is why Eq.~\eqref{eq:dictionary} is needed. After elimination, $\nabla_\mu T_B^{\mu\nu}=0$ rather than a guessed conservation law for a bare canonical tensor is imposed.

Equation \eqref{eq:fulltransport} is equivalent to a projected spin-density transport equation. If $a_\mu=u^\alpha\nabla_\alpha u_\mu$, then
\begin{equation}
 u^\lambda\nabla_\lambda b_{\mu\nu}+\Theta b_{\mu\nu}
 =-2u_{[\mu}b_{\nu]\lambda}a^\lambda.
\end{equation}
Its contraction gives $\Dt b^2=-2\Theta b^2$ in this ideal nondissipative formulation. This is a published transport specialization; collision, relaxation or particle-production laws are not being supplied.

\subsubsection{Exact FLRW versus statistical FLRW}
A nonzero spatial two-form is dual to a nonzero spatial axial vector and selects a direction. Exact FLRW invariance of the spin/torsion configuration under the spatial isometry group forces $b_{\mu\nu}=0$ for this convective Frenkel source. An isotropic axial torsion three-form is a different possible EC source and must not be identified with $u^\lambda b_{\mu\nu}$. The coherent nonzero model therefore has \emph{no exact strong-symmetry FLRW realization}. This agrees with the symmetry warning in \cite{BB}. A statistically isotropic ensemble can have zero mean two-form but nonzero variance and an FLRW \emph{effective metric}; it then requires the averaging assumptions of Model 2 or a separately closed moment theory. There is no contradiction between these two statements of symmetry.

\paragraph{A verified metric-only exception.}
The preceding exclusion concerns the \emph{complete configuration}, not the metric alone. To prevent an overstatement, consider the selected ideal field formulation on a flat FLRW metric in Cartesian coordinates, with a spatially homogeneous aligned spin $b_{xy}=a^2\beta(t)$, all other spatial components zero. Its scalar is $b^2=\beta^2$, and direct substitution gives $\mathcal D_{\mu\nu}=0$ and $\dot\beta+3H\beta=0$. In spherical coordinates the same tensor has
\begin{equation}\label{eq:aligned}
 b_{r\phi}=a^2\beta r\sin^2\theta,\qquad
 b_{\theta\phi}=a^2\beta r^2\sin\theta\cos\theta.
\end{equation}
It is a smooth Cartesian aligned configuration, not a radial-spin profile with a singular center. The metric equations close as Eq.~\eqref{eq:homogeneous} with $q=\kap\beta^2$ and $\dot\beta=-3H\beta$. Thus a nonzero coherent spin field can be compatible with an \emph{FLRW metric} while failing rotational invariance of the full matter/torsion configuration. This is a symmetry-noninheriting source. The cancellation and its nonzero rotational Lie derivative are independently checked. It is an analytical compatibility example, not a proposed astrophysical spin preparation, and no classification of all such branches for $k\ne0$ is claimed. A generic ``full-spin FLRW'' mode would have to state which notion of symmetry and which formulation it uses.

\subsubsection{Exact spherical spin branch}
Under proper rotations about a regular center, a smooth covariant two-form has invariant basis $\dd t\wedge\dd r$ and $\sin\theta\,\dd\theta\wedge\dd\phi$. Frenkel with comoving $u$ kills the first. Thus the coherent ideal branch must have
\begin{equation}\label{eq:radspin}
 b_{\theta\phi}=R^2\sin\theta\,b(t,r),\qquad b_{\phi\theta}=-b_{\theta\phi},\qquad\text{all other components }0.
\end{equation}
This corresponds to a radial axial spin vector. It is allowed by $SO(3)$, but reflection/parity invariance of an axial radial field would force it to vanish. Smoothness at the center requires $b(t,r)=O(r)$ in a regular Cartesian extension, not a constant nonzero radial spin magnitude at $r=0$.

There is a particularly restrictive result. Substitute Eq.~\eqref{eq:radspin} into the retained derivative tensor \eqref{eq:Dspin}. Angular divergence vanishes; radial acceleration contracts with no spin component; angular shear eigenvalues coincide and their symmetric contraction with an antisymmetric angular spin tensor cancels. Hence
\begin{equation}\label{eq:radresult}
 \mathcal D_{\mu\nu}=0,\quad Q=0,\quad P_r=P_\perp=w\rho-\kap b^2,\quad\eps=\rho-\kap b^2,
 \quad\partial_t(BR^2b)=0.
\end{equation}
The cancellation is checked for arbitrary $\Phi,\Lambda,R,b$ by covariant differentiation, as is the Lie invariance under all three rotation generators. A nonspherical spin component produces nonzero derivative stress in a separate control test. The cancellation therefore does not silently discard terms: it follows from this precise symmetry-restricted ideal formulation. Tidal shear alone does not generate a thermodynamic $P_r\ne P_\perp$ here.

Equations \eqref{eq:radresult} plus $p=w\rho$ give a formally closed smooth branch on the same nonsingular-fluid domain as Model 2, with $q=\kap b^2$ and the same local acoustic obstruction. It remains distinct from an unpolarized ensemble. A regular nonzero coherent radial branch cannot be promoted to an exact full-FLRW spin configuration. If the proposed Model 3 instead retains nonzero averaged velocity--spin, shear--spin, spin-gradient or metric correlations, these require additional moment evolution and constitutive equations. Equations for arbitrary $Q,P_r,P_\perp$ in Section 2 are a source decomposition, not that missing closure. This broader proposal remains \status{not closed; implementation blocked}.
